\section{Discussion}

\para{The Curvature Crisis.} The 480$\times$ spike in the information metric reveals that the model's representation space is far from being globally linear. While local linearity holds for semantically related clusters, moving between these clusters requires a ``jump'' over regions of extreme non-linearity. This suggests that the \lrh should be viewed as a local property of semantically coherent regions rather than a global characteristic of the activation space.

\para{Off-Manifold Drift and Robust Steering.} Our findings have direct implications for model steering. Linear interventions, such as those used in CAV-based steering, assume that we can move along a concept vector without exiting the model's learned manifold. However, the 30.7\% increase in \sae reconstruction error for \incongruent paths suggests that forced movement between unrelated concepts drifts into ``unlearned'' regions. This explains why over-steering or steering toward orthogonal concepts often leads to nonsensical model behavior or hallucinations—the model is literally operating in a state it was never trained on.

\para{Unexpected Sparsification.} The observation that forced incongruence leads to a massive decrease in $\normlzero$ sparsity is counter-intuitive. One might expect that a superposition of two unrelated concepts would increase the number of active features. Instead, the model appears to ``collapse'' into a low-dimensional noise state. This suggests that \incongruent inputs may trigger a defensive or default behavior in the model, where specialized semantic features are inhibited.

\para{Limitations.} Our study is limited by the scale of \gpttwo. Larger models, such as Llama-3-70B, might exhibit different geometric properties due to more efficient packing of features or a more robustly learned manifold. Furthermore, the quality of our findings depends on the fidelity of the \saes used for probing; as \sae technology improves, we may gain even clearer insights into the ``off-manifold'' drift.
